\par\Needspace{6\baselineskip}\medskip
\begin{center}\hypertarget{native-187-article140}{}\pdfbookmark[2]{140.}{article-140}140.\end{center}
Ist zunächst $N$ ein beliebiger Punkt des Großkreises $CC'C''$ und zählt man die Abstände der Punkte $C,C',C''$ von $N$ in derselben Richtung, die von $C$ nach $C''$ führt, sodass allgemein gilt
\[
NC''-NC'=2f,\qquad NC''-NC=2f',\qquad NC'-NC=2f'',
\]
so erhält man die Gleichung
\[
\text{I.}\qquad 0=\sin2f\sin NC-\sin2f'\sin NC'+\sin2f''\sin NC''.
\]
Nun wählen wir $N$ im Schnitt der Großkreise $BB^*B''$, $CC'C''$, und zwar als aufsteigenden Knoten des ersten Großkreises auf dem zweiten. Mit $\mathfrak{C},\mathfrak{C}',\mathfrak{C}'',\mathfrak{D},\mathfrak{D}',\mathfrak{D}''$ bezeichnen wir der Reihe nach die Abstände der Punkte $C,C',C'',D,D',D''$ vom Großkreis $BB^*B''$, die auf der einen Seite positiv und auf der gegenüberliegenden Seite negativ genommen werden. Daraus folgt unmittelbar, dass $\sin\mathfrak{C},\sin\mathfrak{C}',\sin\mathfrak{C}''$ jeweils proportional zu $\sin NC,\sin NC',\sin NC''$ sind. Gleichung I nimmt damit die Form an
\[
0=\sin2f\sin\mathfrak{C}-\sin2f'\sin\mathfrak{C}'+\sin2f''\sin\mathfrak{C}''
\]
oder, nach Multiplikation mit $rr'r''$,
\hypertarget{native-188}{}
\[
\text{II.}\qquad 0=nr\sin\mathfrak{C}-n'r'\sin\mathfrak{C}'+n''r''\sin\mathfrak{C}''.
\]
Ferner verhält sich offenbar $\sin\mathfrak{C}$ zu $\sin\mathfrak{D}'$ wie der Sinus des Abstands des Punktes $C$ von $B$ zum Sinus des Abstands des Punktes $D'$ von $B$, wobei beide Abstände in derselben Richtung gemessen werden. Daraus folgt
\[
-\sin\mathfrak{C}=\frac{\sin\mathfrak{D}'\sin CB}{\sin(AD'-\delta)}\mathpunct{\raisebox{.45ex}{,}}
\]
und ganz entsprechend ergibt sich
\[
\begin{aligned}
-\sin\mathfrak{C}&=\frac{\sin\mathfrak{D}''\sin CB}{\sin(AD''-\delta)}\\[.8ex]
-\sin\mathfrak{C}'&=\frac{\sin\mathfrak{D}\sin C'B^*}{\sin(A'D-\delta'+\sigma)}
=\frac{\sin\mathfrak{D}''\sin C'B^*}{\sin(A'D''-\delta'+\sigma)}\\[.8ex]
-\sin\mathfrak{C}''&=\frac{\sin\mathfrak{D}\sin C''B''}{\sin(A''D-\delta'')}
=\frac{\sin\mathfrak{D}'\sin C''B''}{\sin(A''D'-\delta'')}\mathpunct{\cdot}
\end{aligned}
\]
Teilt man also Gleichung II durch $r''\sin\mathfrak{C}''$, so erhält man
\[
\begin{aligned}
0={}&n\cdot\frac{r\sin CB}{r''\sin C''B''}\cdot\frac{\sin A''D'-\delta''}{\sin(AD'-\delta)}\\[.8ex]
&-n'\cdot\frac{r'\sin C'B^*}{r''\sin C''B''}\cdot\frac{\sin(A''D-\delta'')}{\sin(A'D-\delta'+\sigma)}+n''.
\end{aligned}
\]
Bezeichnen wir nun den Bogen $C'B'$ mit $z$, setzen für $r,r',r''$ ihre Werte aus dem vorigen Artikel ein und setzen der Kürze halber
\[
\begin{aligned}
11)\quad&\frac{R\sin\delta\sin(A''D'-\delta'')}{R''\sin\delta''\sin(AD'-\delta)}=a\\[.8ex]
12)\quad&\frac{R'\sin\delta'\sin(A''D-\delta'')}{R''\sin\delta''\sin(A'D-\delta'+\sigma)}=b,
\end{aligned}
\]
so nimmt unsere Gleichung die Form an
\[
\text{III.}\qquad 0=an-bn'\cdot\frac{\sin(z-\sigma)}{\sin z}+n''.
\]
Der Koeffizient $b$ lässt sich auch mit der folgenden Formel berechnen, die sich leicht aus den soeben angegebenen Gleichungen ableiten lässt.
\[
13)\qquad a\cdot\frac{R'\sin\delta'\sin(AD''-\delta)}{R\sin\delta\sin(A'D''-\delta'+\sigma)}=b.
\]
Zur Kontrolle der Rechnung ist es zweckmäßig, beide Formeln 12 und 13 zu verwenden.\HistoricalUnmarkedNote{[Handschriftliche Bemerkung:] Zu Abb. 4: Sei $P$ der Pol des Großkreises $CC'C''$. Dann gilt
\[
\begin{aligned}
&\left(\frac{n''\cos BP}{N''\sin\beta}-1\right)\sin A'D''.\sin B'D\\[.7ex]
&\qquad=\left(\frac{n\cos B''P}{N\sin\beta''}-1\right)\sin A'D.\sin B'D''
\end{aligned}
\]
[wo $N,N''$ die doppelten Dreiecksflächen zwischen den Radiusvektoren der Erde sind.]}
\hypertarget{native-189}{}Ist $\sin(A'D''-\delta'+\sigma)$ größer als $\sin(A'D-\delta'+\sigma)$, so wirken sich die unvermeidlichen Fehler der Tafeln auf die zweite Formel weniger aus als auf die erste. Die zweite ist daher vorzuziehen, wenn die Werte von $b$ eine kleine Abweichung zeigen, die sich daraus erklären lässt. Umgekehrt ist die erste Formel zuverlässiger, wenn $\sin(A'D''-\delta'+\sigma)$ kleiner als $\sin(A'D-\delta'+\sigma)$ ist. Nach Belieben kann man auch einen geeigneten Mittelwert aus beiden Werten wählen.

Zur Überprüfung der Rechnung können auch die folgenden Formeln dienen. Ihre recht einfache Herleitung übergehen wir der Kürze halber.
\[
\begin{aligned}
0&=\frac{a\sin(l''-l')}{R}-\frac{b\sin(l''-l)}{R'}\cdot\frac{\sin(\delta'-\sigma)}{\sin\delta'}+\frac{\sin(l'-l)}{R''}\\[.8ex]
b&=\frac{R'\sin\delta'}{R''\sin\delta''}\cdot\frac{U\cos\beta\cos\beta''}{\sin(AD'-\delta)\sin\epsilon'}\mathpunct{\raisebox{.45ex}{,}}
\end{aligned}
\]
Dabei bezeichnet $U$ den Quotienten $\frac{S}{\sin(\delta'-\sigma)}=\frac{T\sin(t+\gamma')}{\cos(\delta'-\sigma)}$ (Artikel 138, Gleichung 10).
